\documentclass[a4paper]{standalone}
\usepackage{graphicx}
\usepackage{amsmath,amssymb}
\usepackage{hyperref}

\title{Efficient coding hypothesis}
\begin{document}
    \maketitle
    The efficient coding hypothesis was proposed by Horace Barlow in 1961 as a theoretical model of sensory neuroscience in the brain.[1] Within the brain, neurons communicate with one another by sending electrical impulses referred to as action potentials or spikes.
    
    Barlow hypothesized that the spikes in the sensory system formed a neural code for efficiently representing sensory information. By efficient it is understood that the code minimized the number of spikes needed to transmit a given signal. This is somewhat analogous to transmitting information across the internet, where different file formats can be used to transmit a given image. Different file formats require different numbers of bits for representing the same image at a given distortion level, and some are better suited for representing certain classes of images than others. According to this model, the brain is thought to use a code which is suited for representing visual and audio information which is representative of an organism's natural environment .
    \url{https://en.wikipedia.org/wiki/Efficient_coding_hypothesis}
\end{document}